% ============================================================
%  Chapter 5 — Undersea Cables
%  The Secret Life of the Internet
% ============================================================

\chapter{Undersea Cables}

\chapteropening{%
  The internet can send a message from one side of the world to the other in less than a second.
  How?  By racing through giant cables laid on the floor of the ocean!
}
\chapterbadge{Deep-Sea Data Adventure}

% --- Opening story ---
\section*{Crossing the Ocean}

Dot had been sent all the way from Sam's tablet in one country
to a server in another — but to get there, Dot had to cross the sea.

There were no bridges.
No aeroplanes for data packets.

Instead, Dot arrived at a small building by the beach.
The building looked ordinary from the outside, but inside was something extraordinary:
a thick cable, smooth and dark, disappearing into the water.

A gruff but friendly voice echoed up from the deep.

\character{Captain Cable}{Ahoy, Dot!  Welcome to the cable station.
  I am the longest traveller in the internet —
  I stretch from this beach all the way to the other side of the ocean.
  Hop in, and I'll carry you at the speed of light!}

\sectionrule

% --- What are undersea cables? ---
\section*{What Are Undersea Cables?}

\term{Undersea cables} (also called \term{submarine cables}) are thick bundles of cables
laid on the floor of the ocean.
They carry internet data between continents.

These cables are incredibly important.
When you watch a video from another country, send an email to a faraway friend,
or look at a website based overseas — your data almost certainly travels through one of these cables.

\begin{funfact}
The first transatlantic telegraph cable was laid in 1858. Queen Victoria and President Buchanan exchanged the first messages — it took 16 hours to transmit 98 words!
\end{funfact}

\begin{picturethis}
  Imagine a spaghetti noodle stretched from one side of the ocean to the other.
  Now imagine it was actually a bundle of \textbf{optical fibres} —
  threads thinner than a human hair that carry flashes of light.
  Each flash of light is a piece of data.
  Billions of flashes per second race through those thin threads.
  \textbf{That is how an undersea cable works.}
\end{picturethis}

\begin{diagram}[Cross-section of an undersea cable and its protective layers.]
  \begin{tikzpicture}[>=Stealth]
    \def\cx{0}
    \def\cy{0}
    \fill[black!70] (\cx,\cy) circle (1.6);
    \fill[WarmOrange!70] (\cx,\cy) circle (1.3);
    \fill[gray!60] (\cx,\cy) circle (1.0);
    \fill[gray!20] (\cx,\cy) circle (0.72);
    \fill[SkyBlue!45] (\cx,\cy) circle (0.42);

    \draw[white,thick] (\cx,\cy) circle (1.6);
    \draw[white,thick] (\cx,\cy) circle (1.3);
    \draw[white,thick] (\cx,\cy) circle (1.0);
    \draw[white,thick] (\cx,\cy) circle (0.72);
    \draw[white,thick] (\cx,\cy) circle (0.42);

    \node[font=\scriptsize\bfseries,fill=white,inner sep=1pt] at (0,1.48) {1};
    \node[font=\scriptsize\bfseries,fill=white,inner sep=1pt] at (1.18,0.42) {2};
    \node[font=\scriptsize\bfseries,fill=white,inner sep=1pt] at (0.95,-0.48) {3};
    \node[font=\scriptsize\bfseries,fill=white,inner sep=1pt] at (0,-0.82) {4};
    \node[font=\scriptsize\bfseries,fill=white,inner sep=1pt] at (0,-0.28) {5};
  \end{tikzpicture}

  \vspace{6pt}
  {\small
  \begin{tabularx}{0.88\linewidth}{>{\bfseries}c Y}
  1 & Polyethylene layer \\
  2 & Copper power core \\
  3 & Steel strength wires \\
  4 & Gel buffer \\
  5 & Optical fibres \\
  \end{tabularx}
  }
\end{diagram}

\sectionrule

% --- How they are built ---
\section*{How Are Undersea Cables Made and Laid?}

Building an undersea cable is an enormous engineering project.

\begin{quickmap}
  \textbf{Building a submarine cable, step by step:}
  \begin{enumerate}[label=\textbf{\arabic*.}]
    \item Engineers design a route across the ocean floor, avoiding mountains and active volcanoes.
    \item A special ship called a \term{cable-laying ship} carries thousands of kilometres of cable, coiled in its hold.
    \item The ship sails slowly across the ocean, lowering cable carefully over the stern.
    \item Near the shore, the cable is buried under the seabed to protect it from anchors and fishing gear.
    \item In deep water, the cable rests gently on the ocean floor.
    \item Engineers at each end connect the cable to the internet network on land.
  \end{enumerate}
\end{quickmap}

A single cable can be longer than the distance from Earth to the Moon!

\sectionrule

% --- What is inside the cable? ---
\section*{Inside the Cable}

From the outside, an undersea cable looks like a thick rope.
But look inside and you find many carefully engineered layers:

\begin{itemize}
  \item \textbf{Optical fibres} — the core. Hair-thin glass threads that carry data as pulses of light.
  \item \textbf{Protective layers} — plastic, steel wires, and waterproof coatings that shield the fibres from pressure and salt water.
  \item \textbf{Power conductor} — a copper wire that carries electricity to power amplifiers along the cable.
\end{itemize}

\begin{picturethis}
  Think of the cable like a chocolate-coated biscuit.
  The chocolate outside protects the biscuit inside.
  The optical fibre is the delicious biscuit —
  the part that actually does the important job.
  Everything else is protection.
\end{picturethis}

\sectionrule

% --- Scale ---
\section*{How Many Cables Are There?}

There are more than 400 active undersea cable systems around the world,
totalling over 1.3 million kilometres of cable.

If you laid them end to end, they would reach around the Earth more than 30 times!

These cables carry about 99\% of all internet data that travels between continents.
Despite all the talk of satellites, it is really cables under the ocean that keep the world connected.

\sectionrule

% --- Threats and repairs ---
\section*{Can Cables Break?}

Yes, occasionally.
Cables can be damaged by:
\begin{itemize}
  \item Ships' anchors
  \item Fishing gear
  \item Underwater landslides
  \item Shark bites (sharks are occasionally attracted to the electromagnetic fields)
\end{itemize}

When a cable breaks, a special repair ship sails out, hooks the cable up from the ocean floor,
fixes the break, and lowers it back down.
It is like surgery on the ocean floor!

\sectionrule

\begin{newwords}
  \begin{description}[labelwidth=3.8cm, leftmargin=4cm]
    \item[\term{Undersea cable}]    A cable on the ocean floor that carries internet data between continents.
    \item[\term{Optical fibre}]     A thin glass thread that carries data as pulses of light.
    \item[\term{Cable-laying ship}] A ship designed to carry and place undersea cables across the ocean.
    \item[\term{Amplifier}]         A device along the cable that boosts the light signal so it does not fade over distance.
  \end{description}
\end{newwords}

\sectionrule

\begin{tryit}
  \textbf{Draw the ocean route!}

  On a map of the world (or draw your own simple outline), draw:
  \begin{itemize}
    \item A dot on North America — that is Sam's tablet sending a message.
    \item A dot on Europe — that is the server receiving it.
    \item A line along the bottom of the Atlantic Ocean connecting the two dots.
    \item A small ship icon near the top of the ocean (a cable-laying ship!).
    \item Some small fish around the cable on the ocean floor.
  \end{itemize}
  \emph{That glowing line on the ocean floor?  That is Dot's highway!}
\end{tryit}

\sectionrule

\begin{bigidea}
  Most internet data that travels between countries and continents
  moves through \textbf{undersea cables} — bundles of optical fibres on the ocean floor
  that carry information as pulses of light, faster than you can blink.
\end{bigidea}

\character{Captain Cable}{Safe travels, Dot!
  You have crossed the ocean.
  Now, on the other side, let's see how all those tiny messages and videos
  are broken into pieces for the journey.
  Meet the packet crew!}
